\section{Описание модели}
\label{sec:model}

Для каждого изображения обучается отдельная линейная рециркуляционная сеть с архитектурой $n \to p \to n$. Сеть содержит только два линейных слоя и не использует смещения и нелинейные функции активации.

Особенности реализации, соответствующие варианту 2:
\begin{enumerate}
	\item \textbf{Адаптивный коэффициент обучения.} Шаг вычисляется для каждого образа по формуле~\eqref{eq:alpha}; масштаб $\beta$ задаётся пользователем.
	\item \textbf{Ненормированные веса.} После обновления матрицы $W$ и $W'$ не масштабируются и не приводятся к единичной норме.
	\item \textbf{Полное покрытие изображения.} Позиции блоков строятся так, чтобы покрыть всё изображение; при необходимости крайний блок сдвигается к границе.
	\item \textbf{Перекрытие блоков.} Шаги $s_h$ и $s_w$ могут быть меньше размеров блока; при сборке перекрывающиеся значения усредняются.
	\item \textbf{Стохастический порядок.} На каждой эпохе порядок блоков перемешивается генератором с фиксированным seed, что обеспечивает воспроизводимость.
	\item \textbf{Критерий остановки.} Обучение прекращается при $E \le \varepsilon$ или при достижении $T_{\max}$.
\end{enumerate}

Один блок цветного изображения рассматривается как единый вектор длины $n=r\,m\,|S|$; каналы R, G и B записываются последовательно. Все блоки одного изображения проходят через одну и ту же пару матриц $W$, $W'$, поэтому декодер хранится один раз на изображение.

Основные серии экспериментов выполнены при $r=m=8$, $s_h=s_w=8$ (без перекрытия), $p=16$, $\beta=1$, $T_{\max}=80$. Отдельная серия с $s_h=s_w=4$ (перекрытие 50\%) показывает влияние перекрытия на ошибку и коэффициент сжатия.
